\section{Анализ предметной области и обзор существующих подходов}


\subsection{Особенности образовательного видео как объекта оценки}


Образовательное видео отличается от развлекательного по ряду характеристик, принципиальных для задачи оценки.
Во-первых, средняя продолжительность лекционного контента на порядок превышает длину развлекательного ролика,
что усложняет анализ и делает глобальные агрегированные оценки малоинформативными.
Во-вторых, контент мультимодален: речь лектора, слайды, жесты и мимика несут смысловую нагрузку одновременно и по-разному влияют на восприятие в зависимости от аудитории.
В-третьих, интерактивность распределена неравномерно — зрительская активность концентрируется в отдельных точках видео и не отражает качество остального материала.
Наконец, релевантность образовательного видео не является универсальной характеристикой: одна и та же лекция может быть ценной для студента первого курса
и бесполезной для практикующего специалиста. Оценка, не учитывающая контекст пользователя, принципиально неполна.

Из этих особенностей следует, что метрики, разработанные для оценки развлекательного контента (удержание, лайки, просмотры,
CTR) не применимы к образовательному видео: они оптимизированы под другую цель и не несут содержательной информации о качестве подачи или релевантности материала.

\subsection{Существующие инструменты и их ограничения}
Автоматическая оценка образовательного видеоконтента в существующих системах сводится преимущественно к двум категориям метрик:
технические атрибуты (разрешение, частота кадров, уровень аудио) и тематическая классификация (тема, ключевые концепции, наличие нежелательного контента) \cite{zhu2025vceval}.
Эти метрики используются платформами для помощи пользователям при выборе контента, однако не затрагивают ни качество подачи, ни педагогическую структуру.
В академической литературе предпринимались попытки оценить более глубокие аспекты.
\textcite{shi2019lecturevideoquality} исследуют корреляции мультимодальных признаков (аудио, текстовых и визуальных) с субъективными оценками качества лекций на MOOC-платформе,
устанавливая что просодические характеристики и кросс-модальные признаки значимо предсказывают воспринимаемое качество и прирост знаний.
\textcite{pan2022multimodal} предлагают мультимодальный классификатор, обученный предсказывать качество преподавания по пяти аспектам (Clarity, Empathy, Classroom Interaction,
Technical Management и Time Management) на размеченных записях очных и вебинарных занятий, достигая точности $0.898$ на бинарной классификации.
\textcite{zhu2025vceval} предлагает иную постановку: модель симулирует студента, «посещает» лекцию, затем сдаёт тест по внешним учебным материалам темы,
и качество видео приравнивается к результату этого теста.
Объединяет все три подхода базовая установка: задача формулируется как автоматическое вынесение вердикта о качестве видео.
При этом VCEval требует внешних учебных материалов по каждой теме, что ограничивает применимость к произвольному контенту, а
\textcite{shi2019lecturevideoquality} и \textcite{pan2022multimodal} требуют размеченных обучающих данных и плохо переносятся на новые домены и форматы подачи.

Gemini API позволяет взаимодействовать с YouTube-видео напрямую по ссылке и генерировать суммаризации и таймкоды \cite{googledevelopers2024gemini2video}.
Недавно анонсированный Ask YouTube расширяет эту функциональность в сторону поиска по личной истории просмотров \cite{peters2026askyoutube,thekeyword2026askyoutube}.
Оба инструмента позиционируются как средства работы с контентом в режиме диалога, а не как аналитические системы оценки. Стоимость инференса по YouTube URL через Gemini API
может быть существенной при масштабной обработке:
видео токенизируется примерно как $300$ токенов в секунду, поэтому один час видео даёт около $1.08$ млн входных токенов.
В зависимости от модели это составляет примерно от $\$0.11$ за час видео на Gemini 2.5 Flash-Lite до $\$2.70+$ на Gemini 2.5 Pro, без учёта стоимости выходных токенов \cite{googledevelopers2026geminivideounderstanding}.
Это ограничивает применимость подхода при массовом анализе YouTube-видео.

Платформы аналогичные StudyFetch больше ориентированы на образование, но решают задачу усвоения уже выбранного материала: генерация конспектов, флеш-карт,
тестов \cite{studyfetch2024documentation,laxuai2026flashcardcomparison}.
Они включаются после выбора видео и не затрагивают вопрос о том, подходит ли это видео конкретному пользователю.

Поскольку ни один из перечисленных подходов не ставит задачу иначе:
предоставить пользователю структурированную интерпретируемую аналитику по нескольким измерениям одновременно (форме, содержанию и стилю подачи),
чтобы он мог сам принять решение о соответствии видео своим целям. Ниже рассматриваются методологические основы для каждого из них:
методы оценки технического качества аудио и видео (\cref{subsec:technical-av-quality}),
подходы к анализу содержания и тематической сегментации (\cref{subsec:content-analysis-segmentation})
и психологические функции как инструмент формализации паттернов восприятия (\cref{subsec:psychological-functions-flags}).

\subsection{Методы оценки технического аудио-визуального качества}
\label{subsec:technical-av-quality}


\paragraph{Качество аудио.}


Методы оценки качества речевого аудио делятся на intrusive и non-intrusive.
Intrusive-методы сравнивают сигнал с эталоном: классический пример — PESQ (Perceptual Evaluation of Speech Quality) \cite{itut2001pesq},
стандарт для оценки качества телефонной речи. Такие методы дают точные результаты в контролируемых условиях, но неприменимы к пользовательскому контенту,
где эталонная запись отсутствует по определению.

Non-intrusive методы предсказывают воспринимаемое качество непосредственно из сигнала.
Среди них выделяется класс перцептивных моделей, обученных имитировать субъективную оценку слушателя по шкале MOS (Mean Opinion Score).
Одним из актуальных представителей является DNSMOS P.835 \cite{reddy2022dnsmosp835}, предсказывающая три независимых компонента: качество речевого сигнала,
уровень фонового шума и интегральное качество.

Для анализа просодики и характеристик голоса используются инструменты фонетического анализа — в частности, Praat \cite{boersma2001praat},
обеспечивающий точное извлечение фундаментальной частоты (F0) и интенсивности.
Монотонность речи как характеристика подачи устойчиво коррелирует с низкой вариативностью F0 и громкости \cite{shi2019lecturevideoquality}.
Для оценки эмоциональной окраски речи активно развиваются трансформерные модели \cite{wagner2022transformerser}, которые показывают,
что предыдущие подходы систематически занижали точность по валентности, и предлагают архитектуру, закрывающую эту нишу.

\paragraph{Качество видео.}


No-reference оценка качества видео — активная область исследований, особенно применительно к пользовательскому контенту (UGC).
Reference-based метрики — PSNR, SSIM и VMAF \cite{li2016vmaf} — показывают высокую точность в задачах сжатия и стриминга, где оригинальный сигнал доступен,
однако неприменимы к оценке снятого пользователем контента без эталона. VMAF, в частности, активно используется Netflix для контроля качества трансляций,
но требует оригинального мастер-файла для сравнения.

Классические no-reference подходы, такие как BRISQUE \cite{mittal2012brisque}, работают на уровне отдельных кадров и оценивают резкость и артефакты в пространственной области.
Более современные модели, например DOVER \cite{wu2023vqa}, рассматривают видео целиком и разделяют техническую и эстетическую составляющие качества,
что актуально для контента с неоднородными условиями съёмки.

Отдельный класс задач — анализ визуальной сцены и поведения спикера.
Для детектирования и трекинга человека в кадре применяются фреймворки компьютерного зрения реального времени,
среди которых MediaPipe \cite{lugaresi2019mediapipe} обеспечивает надёжное извлечение ключевых точек лица и позы.
Обнаружение монтажных переходов решается через детекторы смены сцен \cite{castellano2020pyscenedetect}.
Задача распознавания текста в кадре — на доске или слайдах — решается средствами OCR.

\subsection{Анализ содержания и тематическая сегментация}
\label{subsec:content-analysis-segmentation}

Задача тематической сегментации длинного текста или транскрипта — разбиение на связные смысловые блоки — изучается давно.
Классический подход TextTiling \cite{hearst1997texttiling} основан на лексической когезии:
границы тем определяются по падению лексического сходства между соседними окнами текста. Метод интерпретируем и не требует обучения,
однако чувствителен к шуму в транскрипте и плохо справляется с плавными тематическими переходами, характерными для лекционного контента.

Нейросетевые подходы на основе эмбеддингов улучшают качество сегментации за счёт семантического представления текста,
но по-прежнему работают с транскриптом как единственным модальным каналом.
LLM-based подходы — в частности, Chapter-Llama \cite{ventura2025chapterllama} — используют языковые модели для одновременной сегментации и именования глав,
приближаясь по качеству к авторской разметке на видео с чёткой структурой.

\subsection{Психологические функции и флаги восприятия}
\label{subsec:psychological-functions-flags}

Типология Юнга \cite{jung1921psychologischetypen} описывает четыре базовые функции обработки информации:
мышление (Логик), ощущение (Сенсорик), интуиция (Интуит) и чувство (Эмоция).
Типология получила широкое практическое применение через адаптацию MBTI \cite{myers1998mbtimanual} — прежде всего в HR и карьерном консультировании.

Для задачи оценки образовательного видео интерес представляет не личностная интерпретация типологии,
а описание каналов восприятия: каждая функция определяет, через какой модальный канал человек преимущественно воспринимает информацию.
Логический тип ориентируется прежде всего на аудиальный канал — звук, чёткость речи, — а также на фактологическую структуру и системность изложения.
Сенсорный — на визуальную конкретику и реальность происходящего в кадре. Интуитивный — на идеи, абстракции и концептуальные связи.
Эмоциональный — на атмосферу, присутствие человека и характер взаимодействия с аудиторией.

\subsection{Выводы}


Анализ существующих подходов показывает, что задача оценки образовательного видеоконтента не решается ни одним из них в отдельности.
Методы оценки технического качества не учитывают содержание; инструменты анализа содержания не принимают во внимание пользовательский контекст;
существующие платформы решают смежные задачи усвоения, а не выбора контента. Отдельную проблему представляет отсутствие интерпретируемости:
большинство подходов выдают единый балл без декомпозиции и объяснения.
Совокупность этих ограничений обосновывает необходимость каскадного мультимодального подхода с явным учётом профиля пользователя.